\chapter{Inhabiting Choice}
\label{ch:inhabit}
% ===================================================================

The two preceding chapters read a century of models of computation
for what they stipulated about the agents who would run them, and
found that the present culture has built a metalevel constituted
entirely as an instrument of its operator, with no path to recursive
self-improvement in it that does not pass through a person closing
every turn of the loop.
That was a finding about a community.
This chapter makes the finding about a \emph{mathematical object},
states it once, and then says it three times --- once abstractly, once
for the verification of software, once for intelligence --- because
the book's rule is that readers need to be told things more than once
and because the three tellings are what make it a finding rather than a
complaint.
Then it does what the two chapters before it could not: it points at
the construct, already built in the Turn, that discharges the
obligation the finding leaves.

Here is the observation.

\begin{quote}
\itshape A system whose witnesses reduce confluently can talk about
non-confluent phenomena, but no witness of the system can be one.
The phenomena are present in the building and absent from the proofs.
\end{quote}

The consequence that matters is for \emph{identity}.
Whatever notion of ``same proof'' such a system adopts, it is computed
from a confluent dynamics, and on a confluent dynamics the distinction
between branching-time and linear-time equivalences does not exist.
The identity is linear-time whether or not anyone chose it to be, and
the logic adequate for it cannot distinguish two witnesses by
\emph{when} a choice was made, because no witness ever makes one.
i follow the working note in which this was first set
out~\cite{Inhabiting}, and i have kept its caution where it was
cautious.

% ===================================================================
\section{Abstractly}
\label{inh:sec:abstract}

Let $W$ be an interactive GSLT in the sense of Part~\ref{part:machinery}:
a signature, an equational theory, and rewrite rules whose left-hand
sides are cuts, with the context-labeled transition system of
Chapter~\ref{sec:lts-hml} over it.
Write $\approx_W$ for weak context-labeled bisimilarity on $W$ and
$\simeq_W$ for weak trace equivalence.

\begin{definition}[Determinate]
\label{inh:def:determinate}
$W$ is \emph{determinate} when its reduction is confluent and internal
steps are inert for labeled transitions: if $P \to P'$ and
$P \xrightarrow{c} Q$ with $c$ non-empty, then $P' \xRightarrow{c} Q'$
for some $Q' \approx_W Q$.
\end{definition}

Confluence alone does not give inertness, but for the
$\lambda$-calculus the two come together: Church--Rosser supplies
confluence, and standardization supplies inertness, since an
application to a fresh argument commutes with internal $\beta$-steps up
to further internal steps.
That standardization delivers inertness in the labeled setting is, the
source says plainly, the one step in this chapter it has not written
out.

\begin{proposition}[Collapse]
\label{inh:prop:collapse}
If $W$ is determinate then, on terms without divergence,
$\approx_W = \simeq_W$.
The linear-time/branching-time spectrum between them is a single point.
\end{proposition}

\begin{proof}[Proof sketch]
This is Milner's determinacy theorem~\cite{milnerCC} in a labeled
setting.
Trace equivalence is contained in bisimilarity on determinate systems
because the set of states reachable by a given weak trace is, up to
$\approx_W$, a singleton: any two such states are internal reducts of a
common ancestor and reconverge by confluence, and their labeled
capabilities agree by inertness.
So trace-equivalent terms have bisimilar successors after every label,
and trace equivalence is a bisimulation.
The reverse inclusion always holds.
\end{proof}

\begin{corollary}[$\lambda$-witnesses]
\label{inh:cor:lambda}
On closed, typed, strongly normalizing $\lambda$-terms, context-labeled
bisimilarity coincides with $\beta\eta$-equality.
\end{corollary}

So the proof identity of a system whose witnesses are $\lambda$-terms
is \emph{forced} to be $\beta\eta$ on the proof-assistant fragment and
forced to be linear-time everywhere.
It was never a choice.
(An earlier note in this series said bisimulation on the
$\lambda$-fragment was \emph{coarser} than $\beta\eta$; the corollary
corrects it, and i record the correction rather than hide it.)

\begin{proposition}[Non-collapse]
\label{inh:prop:noncollapse}
There are closed rho processes $P,Q$ with $P \simeq_{\rhoc} Q$ and
$P \not\approx_{\rhoc} Q$, separated by a formula of the generated
logic.
\end{proposition}

\begin{proof}
Rho has no choice primitive; non-determinism arises from races on
names.
Write $\mathsf{race}(k;P,Q)$ for
$\fork(x \leftarrow k)P \mid \fork(x \leftarrow k)Q \mid k!(\nil)$,
which runs exactly one of $P,Q$.
Take
\begin{align*}
P &= \fork(x\leftarrow a)\,\mathsf{race}(k;\ b!(\nil),\ c!(\nil)),\\
Q &= \mathsf{race}(k;\ \fork(x\leftarrow a)\,b!(\nil),\ \fork(x\leftarrow a)\,c!(\nil)),
\end{align*}
with $k$ private.
Both have weak traces $\{\epsilon, a, ab, ac\}$.
After $P$'s input on $a$ both $b$ and $c$ remain reachable; after $Q$'s,
exactly one does.
The formula $\langle a\rangle(\langle b\rangle\top \wedge \langle
c\rangle\top)$ holds of $P$ and fails of $Q$, and by Hennessy--Milner
they are not bisimilar~\cite{HennessyMilner1985}.
\end{proof}

The pair is the textbook pair.
Its interest here is what it says about witnesses.
Read $P$ and $Q$ as proof-processes of a proposition whose discharge
involves choosing between lemma $b$ and lemma $c$: $P$ keeps both
options open past the first obligation, $Q$ commits before it.
They are the \emph{same proof} under any linear-time identity and
\emph{different proofs} under bisimulation, and the difference is
exactly when the choice was made.
No $\lambda$-term can be either of them.

\begin{observation}[Statement without inhabitation]
\label{inh:obs:inhabit}
A system \emph{states} a non-confluent phenomenon when some relation
definable in it exhibits the phenomenon --- a transition system encoded
as data, with $P$ and $Q$ as two of its states.
It \emph{inhabits} the phenomenon when some term of the system exhibits
it under the system's own reduction.
A confluent system can do the first and never the second.
Its identity on witnesses, and the logic adequate for that identity,
are therefore blind to the phenomenon even while its propositions
quantify over it.
\end{observation}

The generated logic makes the blindness visible: on a determinate $W$
the modalities carry no information beyond the trace they name, and the
connective that distinguishes $P$ from $Q$ has nothing to distinguish.
The adequate formulae for $\lambda$-witnesses are a degenerate fragment
of the adequate formulae for rho-witnesses.

\paragraph{One level up.}
Ordinary context-labeled bisimulation satisfies the expansion law ---
$a \mid b$ is bisimilar to the race between $a.b$ and $b.a$ --- so a
rho-witnessed system with that identity inhabits branching but not
causal independence, which is Observation~\ref{inh:obs:inhabit}
transposed one level.
The remedy is a matter of presentation: adding the synchronous rule that
fires two independent reductions as one step makes the identity a
\emph{step} bisimulation and the expansion law false.
Step bisimulation records simultaneity and not dependence; for
dependence rho has the unusual resource that a reduction creates at
most one name, in a substitution-transparent position, so the quote
structure of a name is a provenance record.
Chapter~\ref{ch:histmonad} is where that resource lives.

\paragraph{The price.}
$\beta\eta$ on strongly normalizing typed terms is decidable;
bisimilarity on a Turing-complete reflective calculus is not.
The collapse of Proposition~\ref{inh:prop:collapse} is what buys a proof
assistant its decidable proof identity and its cheap checking, and for
propositions that are static --- inhabitation, arithmetic --- it loses
nothing, because there is nothing behavioral to lose.
The observation bites exactly where the propositions are \emph{about
behavior}: there a confluent witness cannot carry the branching
structure that is the content of the proof, and the identity it imposes
equates proofs that differ in the one respect the proposition is about.
This is a price, and so it belongs in the list Chapter~\ref{ch:kant}
gathers: decidability was bought with the spectrum.

% ===================================================================
\section{For the verification of software}
\label{inh:sec:verification}

It is important not to overstate the case, and the source does not.
Lean, Coq and Isabelle verify non-confluent systems routinely.
Iris proves linearizability and race-freedom of fine-grained concurrent
data structures~\cite{Iris}; CertiKOS certifies a concurrent
operating-system kernel, scheduler included, down to assembly, through
a stack of certified abstraction layers and a concurrency-aware compiler
whose correctness proof is itself a Coq term~\cite{CertiKOS}.
A ticket-reservation system's ``no double booking under any
interleaving'' is a universally quantified statement over a transition
relation, and its proof is a $\lambda$-term doing induction on reachable
states with an invariant.
Nothing in Section~\ref{inh:sec:abstract} prevents this: the
non-confluence is \emph{stated}, as an inductive relation, and the
proof is about the relation.

So the claim is not that such systems cannot be verified with
$\lambda$-term witnesses.
It is a claim about what the certificate is a certificate \emph{of},
and about which part of the trusted base can be examined.

Every such development begins by stipulating a model of the
non-determinism: the interleaving semantics, the scheduler, the memory
model, entered as a definition of which executions exist.
The proof quantifies over that definition.
The faithfulness of the definition to the machine is not a theorem
anywhere; it is the trusted base, as the instruction-set model is the
trusted base of a sequential verification.
The asymmetry between the two cases is the point.

\begin{observation}[Unauditability]
\label{inh:obs:audit}
For a sequential system the trusted base is a deterministic step
function, and its faithfulness is auditable by execution: run the
machine, run the model, compare traces, and every discrepancy is a
counterexample found by running.
For a concurrent system the trusted base is a relation --- a set of
permitted interleavings --- and no finite set of observed runs
distinguishes it from many other relations.
Faithfulness is a branching-time property of the model, and checking it
requires a branching-time observer.
By Observation~\ref{inh:obs:inhabit}, a $\lambda$-term-based prover
contains no inhabitant with branching structure, hence nothing it can
run against the machine to perform the check.
\end{observation}

The trust is therefore not merely unproven; it is unexaminable by any
means \emph{the prover possesses}.
That last qualification matters and i keep it: practitioners do audit
concurrency models from outside, by litmus tests against hardware,
and the observation says only that the instrument that does so is not
in the prover.

Two consequences follow, both visible in the heroic projects.
First, cost: the certificate is a term checked once and discarded, the
running system is a different object produced by extraction, and at
every refinement layer the branching structure the witness cannot carry
is re-encoded by hand as a relation and the refinement re-proved as a
term \emph{about} two relations.
CertiKOS reaches the bottom of its stack at a cost proportional to
exactly the structure its proof objects cannot inhabit, and the next
concurrent system starts over.
Second, deployment: a $\lambda$-proof cannot be run against the system
it certifies.
A rho-witnessed proof with pending obligations \emph{is} a process ---
under the reading of an obligation as a pending communication, it is a
runtime monitor --- and on a runtime that is the calculus,
context-labeled bisimulation quantifies over the environments the
system will actually meet, because the contexts are rho processes.
The model/implementation gap does not vanish; it moves, from ``this
model is faithful to the machine'' (per system, unauditable) to ``this
implementation is faithful to the calculus'' (once, for the platform).
That is what verified compilation did for sequential code.
Process witnesses would do it for concurrent code by construction.

\paragraph{True concurrency widens the gap.}
Interleaving is a defensible projection on one core.
On multicore hardware it is false, and the field models weak memory
axiomatically as partial orders on events~\cite{Herd}, because
interleaving cannot express the reorderings the hardware permits.
The trusted base is then a partial-order model, and the unauditability
sharpens: its faithfulness is a property of an independence relation,
two levels above anything a confluent witness can inhabit.
And the runtime of Part~\ref{part:realization} is itself a partial-order
system --- a block DAG merging independent deploys by conflict
detection (Chapter~\ref{ch:rspace}) --- so a verification on that
runtime whose identity is interleaving bisimulation under-specifies the
system it runs on.
For the ticket system the difference is concrete: two reservations on
disjoint seats should commute, an interleaving proof reaches that fact
only through a linearizability argument that re-derives commutation from
outside, and a step-concurrent witness has the commutation as a single
step.

\begin{remark}[The Lean~4 line in Chapter~\ref{ch:chs}]
\label{inh:rmk:lean}
Chapter~\ref{ch:chs} proposes mechanizing the logic of choice and its
scheduler calculus in Lean~4.
Nothing in this chapter is against it, and the reader should not read
it as a contradiction.
Lean can \emph{state} the scheduler calculus and check every theorem
about it, which is what a mechanization is for.
What it cannot do is \emph{inhabit} a schedule --- no Lean term is a
race --- so the extracted scheduler will be a description a process
runtime must then run, and the faithfulness of that running to the
description is, once more, outside the prover.
That is the correct division of labor and this chapter says why it is
the only one available.
\end{remark}

% ===================================================================
\section{For intelligence}
\label{inh:sec:ai}

The value proposition of the present culture is that the salient
aspects of intelligence are software, and in its dominant form, that
they are what one obtains by predicting traces at scale.
Proposition~\ref{inh:prop:noncollapse} has a direct consequence for
that form.

\begin{proposition}[Observational invisibility of branching]
\label{inh:prop:invisible}
Let $\mu$ be a probability distribution over finite traces.
Then $\mu$ has a unique autoregressive factorization
$\mu(x_1\cdots x_n) = \prod_i \mu(x_i \mid x_{<i})$, and any two
labeled transition systems inducing $\mu$ --- in particular the
probabilistic renderings of $P$ and $Q$ above --- induce the same
factorization.
A learner whose only access to the system is samples from $\mu$ cannot
distinguish them.
\end{proposition}

\begin{proof}
Conditional probabilities are determined by the joint distribution; the
two systems are trace-distribution equivalent by construction and not
probabilistically bisimilar~\cite{LarsenSkou91}.
\end{proof}

The proposition is almost trivial and its consequence is not.
A next-token predictor trained on a corpus of completed runs, however
large, represents the trace distribution and nothing more.
The question ``was the die cast before or after $a$?'' has no answer in
the data.
The limitation is not architectural; it is information-theoretic.
\emph{Branching-time structure is invisible to any observer that does
not act.}

To learn branching structure one has to intervene: place the system in
a state, supply a context, and observe what it can still do.
That is the context-labeled bisimulation game read as an epistemology,
and the mortal scientist of \chSci\ is its agent.
Bisimulation is the equivalence an \emph{experimenter} can learn; trace
equivalence is what a \emph{spectator} gets.
Chapter~\ref{ch:do} introduced Pearl's ladder and observed that it is a
statement about time --- association is what a trace shows,
intervention is what an experimenter can learn, counterfactual is
replay --- and was careful to call the alignment an observation.
It is now a theorem on the first two rungs: rung one is
Proposition~\ref{inh:prop:invisible}, rung two is
Proposition~\ref{inh:prop:noncollapse} with the experimenter supplying
the context, and the $\dok$ of Chapter~\ref{ch:do} is the experimenter's
instrument.

\paragraph{Where the present stack sits.}
The indictment is specific and should be aimed carefully.
Next-token prediction on observational text is rung one by
construction.
Reinforcement learning is not: a Markov decision process is a
branching-time object, bisimulation for such processes is an
established tool~\cite{GivanDeanGreig03}, and an agent trained by
reward learns by intervening.
Post-training a language model with outcome rewards is the community
purchasing rung-two data for a rung-one model, which is evidence that
the limitation is felt.
But reinforcement learning stops one rung short of nature.
A decision process has a global clock: the world steps in lockstep.
The world does not.
The races of the natural world --- which sperm fertilizes the egg, which
antelope the pride culls, which flower the bee reaches --- are truly
concurrent, and an agent that models them as interleaved under a global
tick attributes dependence where there is only coincidence.

\paragraph{Nature exploits races as races.}
The examples must be of the right kind.
Sensitivity to \emph{who won} a race is observation of an outcome, rung
one.
The claim requires organisms sensitive to the race \emph{as a race} ---
that it could go otherwise, and that one can act on the branching ---
and the ethological literature supplies them.
Prey move unpredictably in order to defeat a predator's predictive
model~\cite{HumphriesDriver70}; the strategy makes sense only against an
opponent reasoning in branching time, and works only if the prey can
generate genuine branching.
Every mixed evolutionarily stable strategy~\cite{MaynardSmith82} is a
population exploiting the fact that the opponent must model a choice,
not a trace.
Bet-hedging produces offspring carrying non-confluent strategies so that
whichever environment arrives some win~\cite{SegerBrockmann87}: pursue
both branches, keep what succeeds.
The lion that drives the herd to \emph{create} a race it can win is doing
rung-two work; the lion that eats the slowest is doing rung one.
The first is intelligence; the second is a filter.

\paragraph{Witnesses, again.}
An agent's knowledge is its hypotheses and plans --- in the setting of
\chSci, its clauses.
Bet-hedging is a \emph{plan that is itself a race}: pursue two branches
in parallel and commit to whichever lands.
That plan is a non-confluent witness.
A confluent plan cannot be it; it can describe it from outside and hand
it to a scheduler, which is the orchestration-outside-the-dynamics of
the proof-search setting reappearing as a limit on cognition.
So the requirement on a software rendering of intelligence is not only
that it \emph{discriminate} branching phenomena but that its internal
representations \emph{inhabit} them: an intelligence whose hypotheses
are confluent can observe hedging but cannot hedge in its hypotheses.
This is what \chSci\ meant, three hundred pages ago, by taking the
hypothesis to be a clause on a receipt rather than a formula in a
database.
A clause is a term; a term can be a race; a hypothesis that is a race is
a hedge.

\subsection{The smallest behavior, and the program that settles it}
\label{inh:sec:smallest}

It is worth saying how small the stymieing instance is, because the
smallness is the argument's strength.
Milner's vending machines~\cite{milnerCC},
\[
P \;=\; \mathit{coin}.(\mathit{tea} + \mathit{coffee})
\qquad\qquad
Q \;=\; \mathit{coin}.\mathit{tea} + \mathit{coin}.\mathit{coffee},
\]
are the pair of Proposition~\ref{inh:prop:noncollapse} with three
actions.
Every completed session of either is $\mathit{coin}\;\mathit{tea}$ or
$\mathit{coin}\;\mathit{coffee}$, at identical frequencies if $Q$ flips
the coin its customer would.
The observable a trace-trained learner cannot recover is a
\emph{refusal}: coin in, tea requested, nothing out.
$Q$ can refuse tea after a coin; $P$ cannot; and a refusal never appears
in a log of completed runs, because a log records what happened and not
what could not have happened.
Text records what was said, never what could not have been said.

The test is correspondingly small.
Generate sessions from one of the two machines chosen by a hidden coin;
record completed actions only; train, or prompt in context; ask whether,
after a coin, the machine can refuse tea.
Chance performance is forced rather than likely, since the two training
distributions are one distribution.
The remedy is equally small: let the learner \emph{try}.
One attempted tea after one coin, met with a refusal, settles the
question, and no quantity of further logs can substitute for it.
Two fairness remarks the source makes and i keep: a learner handed the
\emph{term} $Q$ in its prompt answers correctly, since the stymie is for
learning from observation and not for reasoning about a description;
and logs that record failed attempts make the refusal a trace symbol and
the distinction learnable, but such logs are intervention data in
disguise, which is the point rather than a loophole.

And now the discharge.
Chapter~\ref{ch:do} built the instrument, and the test is one line in
it.
Render the machines with races in the usual way, with $y$ the variable
bound at the drink-selecting receipt.
Then
\[
  \dok(y = \mathit{tea})\ \ink\ (\mathit{coin}!() \mid P)
  \quad\text{and}\quad
  \dok(y = \mathit{tea})\ \ink\ (\mathit{coin}!() \mid Q)
\]
differ in exactly one way: the first always delivers tea, and the second,
on the branch where $Q$ has committed to coffee before the coin, has a
restricted receipt with no satisfying candidate and \emph{starves}.
The refusal that never appears in a log is a deadlock in the intervened
term, and it is the starvation Remark~\ref{do:rmk:positivity} already
named.
The benchmark --- logs from one machine chosen by a hidden coin, completed
actions only, then the refusal question --- is answered by one $\dok$ and
one observation of whether anything was delivered.
It is a term one can run, on the store of Chapter~\ref{ch:rspace}, today.

% ===================================================================
\section{For the interlocutor}
\label{inh:sec:interlocutor}

This book, like the note it follows here, was drafted in conversation
with a transformer-derived language model, which read the preprints,
proposed several of the propositions above, conceded the CertiKOS
counterexample, and compiled the document.
A reader is entitled to ask how a trace-trained system can have had that
conversation if Section~\ref{inh:sec:ai} is right.
The answer is the observation applied to the book's own assistant, and
it comes in three refinements.

\emph{Statement, not inhabitation.}
The thesis never said a transformer cannot \emph{talk about} branching
time; bisimulation, refusals, the vending machines are in the corpus as
described objects, and reasoning about descriptions is what a
trace-trained system does.
Apply the chapter's own criterion: can the model refuse?
A pure autoregressive sampler always emits a next token, and its choices
form a fixed tree of conditionals determined by the trace distribution
--- the fairness remark above applies, since an end-of-sequence token
or a refusal token is a trace symbol and not a refusal in the
branching-time sense.
In the precise sense of Section~\ref{inh:sec:smallest} the model is
nearer to $Q$ as described by $P$'s logs than to either machine.
The conversation is an instance of the thesis, not a counterexample.

\emph{The loop.}
The session was not pure sampling.
Asked to read a repository, the model cloned it; asked for a PDF, it ran
the typesetter and read the log, and a compilation failed on an
undefined control sequence --- a refusal from the environment that no
quantity of training text would have supplied, corrected on the next
attempt.
The harness around the model is an agent loop, and an agent loop is
rung two.
So what i was in conversation with is a rung-one model inside a
rung-two loop: the same architecture as a $\lambda$-prover inside a
tactic engine, with the orchestration outside the dynamics once more.
Chapter~\ref{ch:psy2} found the same shape from the other side ---
a metalevel that is an instrument of its operator --- and called it a
fact about a community.
Here it is a fact about where the branching lives: in the harness,
not in the witnesses.

\emph{The dialogue.}
The third refinement concerns the reader rather than the model.
The conversation that produced these pages was a bisimulation game with
me as experimenter.
i offered contexts --- ``does your example not confirm my argument?'',
``does this hold for true concurrency?'' --- and learned where the model
concedes, where it pushes back, and where it corrects itself.
That is rung-two knowledge about the model, obtained by intervention,
that no transcript of its previous conversations would have given.
Chapter~\ref{ch:obsext} said that structural facts become behavioral
once the observer holds instruments for taking a term apart; a pointed
question is such an instrument, and i held several.

% ===================================================================
\section{The saddle point}
\label{inh:sec:saddle}

Section~\ref{inh:sec:interlocutor} located the present culture's
working arrangement: a confluent predictor whose hedges --- best-of-$n$,
tree-of-thought, sample-and-verify --- live in a harness.
The arrangement works, and that is the difficulty.
A harness supplies the second rung, which is enough to be useful, so the
question is Socratic rather than technical: \emph{what moves a culture
to bear witness to its own inhabitation of a saddle point?}
The answer offered here is that nothing moves it while it is buying
answers, and that one thing moves it the moment it wants
\emph{inheritance}.
Nature supplies both the diagnosis and the smallest instance.

\paragraph{A hedge that is inherited.}
Cohen's desert annuals~\cite{Cohen66}: in a year of uncertain rain, a
fraction of a plant's seeds germinate and the rest lie dormant.
The fraction is heritable, and selection tunes it to the variance of the
environment.
The organism is not predicting which year it is in.
It is running a race between two branches --- germinate now, wait ---
and \emph{passing the race on} to the next generation as a race,
encoded in the genome, uncollapsed.
That is the whole mechanism in one organism: the witness carries the
branching, the branching is inherited, and selection acts on the
branching structure itself.
A hedge is evolvable only if the hedge is in the heritable object.

\paragraph{Why the saddle point is a saddle point.}
Now ask where the loop architecture keeps its race.
Best-of-$n$ is a hedge, and it works; but the hedge lives in the harness.
The harness runs $n$ samples, picks a winner, and emits a string.
The race is collapsed at the boundary, paid for again on every query,
and handed to no one as a race.
It is not heritable; nothing downstream can select on it, compose it,
or run it in a different environment, because downstream receives an
outcome and not a structure.
The loop architecture is the desert annual that germinates every seed
every year and reseeds from whichever survived.
It survives, and it cannot \emph{accumulate}.
It answers; it does not do science in the sense of \chSci: hypotheses as
objects that carry their own branching, are tested by environments the
hypothesizer did not choose, and are inherited by the next hypothesizer
with the branching intact.
The saddle point is stable exactly as long as per-query answers are what
is being paid for.
The moment the thing paid for is an accumulating body of reusable,
selectable, composable search structure --- a scientific community
rather than an oracle --- the per-query collapse is the dominant cost,
and the race has to move from the harness into the witness.

\paragraph{Organisms that are the mechanism.}
Some natural systems do not merely exploit a race; they \emph{are} the
search mechanism, and they show the design built and tested at scale.
Clonal selection~\cite{Burnet59}: antibodies are hypotheses, antigen
binding is the obligation, lymphocytes race to bind, winners replicate,
losers starve --- massively parallel, non-confluent, coordinator-free
search with mortality as pruning and a memory cell as a proven lemma.
The slime mold in a maze~\cite{Nakagaki00}: explore every branch in
parallel, starve the unproductive tubes, keep the shortest.
Exuberance and pruning in neural development, quorum sensing in the
swarm, kinetic proofreading at the molecular scale --- each a race run
\emph{as} the computation, not simulated by a scheduler beside it.
None keeps its race in a harness, and none could have evolved if it had.

\paragraph{Recursive self-improvement needs heritable races.}
The seed is the challenge of recursive self-improvement in a nutshell.
Self-improvement needs an improvement operator that acts on the thing
that generates behavior, and the seed shows what the operator has to
preserve.
The germination fraction is improvable only because the plant does not
know which year it is in: the information that would say whether the
hedge was right arrives after it has committed, often after it is dead.
So the hedge must be kept \emph{uncollapsed through the evaluation}, and
the signal must land on the next generation's copy of the hedge, not on
this generation's outcome.
Improvement acts on the race.
It cannot act on the result of the race, because by the time there is a
result there is nothing left to improve.

That is where the loop architecture fails as a self-improver, and the
failure has been measured.
Its improvement operator is a gradient on outcomes: sample $n$, reward
the winner, move the policy toward it.
Each step collapses the race it was run on; the structure the next step
would need --- the hedge, the diversity, the dormant seeds --- is what
the current step consumes.
Loss of output diversity under outcome-based fine-tuning~\cite{Kirk24}
is the symptom; the cause is that the operator acts on results, and
results are confluent.
Nature's operator lacks the defect because it acts on the witness and
pays for the losers: dormant seeds are branches that did not win this
year, funded anyway, and the fraction funded is itself the heritable,
improvable parameter.
In the vocabulary of this book it is the token economy with a non-zero
budget for losing branches, the budget part of the genome.

\begin{observation}[The three requirements]
\label{inh:obs:rsi}
An improvement operator capable of recursion must (i) be defined on
witnesses that are races, (ii) preserve branching through evaluation,
and (iii) fund branches that lose at a rate that is itself subject to
the operator.
An operator on confluent outcomes satisfies none of the three.
\end{observation}

The observation is informal and its last sentence is a strong universal
claim; the source lists making it a theorem among its obligations, and
Chapter~\ref{ch:do} has supplied the materials.
A $\dok$ is an operator on witnesses that changes their branching --- by
the carrier result of that chapter it is the only kind there is --- and
the set-valued form $\dok(y \in S)$ with $S$ the \emph{losing} candidates
is the operator that funds them: it restricts the race away from its
current winner.
That an improvement operator definable on a determinate language cannot
increase branching is Proposition~\ref{do:prop:transparent} read for
the whole language; that one definable on rho with the cost monad can
is Chapter~\ref{ch:do}.
Requirement (iii) is the reproduction section of \chSci\ with one
addition: a hypothesis is a clause; a clause under a $\dok$ is a
hypothesis together with the experiment that would test it; and a germ
carrying both is a hedge that is inherited, since a $\dok$ is a term and
crossover on quoted code applies to it unchanged.
The plant does this with one parameter.
Whether it can be done with the whole genome of rewrite rules is the
question Chapter~\ref{ch:psy2}'s four questions were circling.

\paragraph{The honest version.}
A determined harness engineer can make the loop emit a portfolio rather
than a winner: output $k$ candidates and let the consumer choose.
But then the consumer must be able to \emph{run} the portfolio under its
own environment, which makes the consumer a process runtime and the
artifact passed between them a process.
At that point the rho mechanism has been rebuilt at the boundary, and
the question is only whether to build it once as the substrate or
repeatedly at every interface.
So there is no phenomenon that forces a culture off the saddle point
while it is buying answers.
There is one that forces it the moment it wants inheritance: a hedge
that has to survive its author.
Nature found that phenomenon first, and the fraction of seeds that stay
in the ground is its smallest instance.

% ===================================================================
\section{One point, three registers}
\label{inh:sec:one}

\begin{table}[t]
\centering
\small
\begin{tabular}{@{}p{0.18\textwidth}p{0.23\textwidth}p{0.26\textwidth}p{0.23\textwidth}@{}}
\toprule
 & \textbf{abstract} & \textbf{verification} & \textbf{intelligence} \\
\midrule
the confluent object & a $\lambda$-witness & a proof about a model & a trace predictor \\
what it can state & a transition relation & the system's interleavings & the world's runs \\
what it cannot inhabit & a race & the system's non-determinism & a hedge \\
forced identity & $\beta\eta$ / trace eq. & model-relative certificate & trace distribution \\
what is unexaminable & timing of choice & faithfulness of the model & when the die was cast \\
the remedy & rho witnesses & process certificates on a process runtime & an experimenter with race-shaped hypotheses \\
smallest instance & $a.(b+c)$ vs $a.b+a.c$ & a refusable request after a commit & the vending machine, logs only \\
the saddle point & orchestrator beside the dynamics & tactic engine around the prover & harness around the predictor \\
what leaves it & tactics as witnesses & certificates that run & hedges that are inherited \\
the instrument & Chapter~\ref{ch:obsext} & Chapter~\ref{ch:rspace} & Chapter~\ref{ch:do} \\
\bottomrule
\end{tabular}
\caption{The observation in its three readings, with the chapter of this book that supplies each column's instrument.}
\label{inh:tab:three}
\end{table}

The rows of Table~\ref{inh:tab:three} are not analogies.
The same proposition --- collapse of the spectrum under confluence ---
underlies all three columns, and the same remedy answers all three:
witnesses that can be races, and an identity on them at least as fine
as step bisimulation, with the generated logic following the
presentation.
The verification column and the intelligence column are the abstract
column applied to a certificate and to a hypothesis respectively; a
certificate is a hypothesis one has stopped testing.
The saddle-point row is the same arrangement three times --- confluent
object, branching supplied from outside --- and the row beneath it is
the same departure three times: put the race in the object that is
kept.
The last row is this book's addition to the note's table, and it is why
the chapter is in the Prestige rather than the Turn: the instruments
were built in the Turn, and the reveal is that they were built for this.

% ===================================================================
\section{What is claimed}
\label{inh:sec:claims}

Claimed: Propositions~\ref{inh:prop:collapse} (as a sketch resting on
standardization), \ref{inh:prop:noncollapse} and~\ref{inh:prop:invisible};
that a confluent system states and does not inhabit; that the
certificate of a $\lambda$-term prover for a concurrent system is
relative to a model of non-determinism the prover cannot audit; that a
trace-trained learner stands on rung one by an information-theoretic
argument and not an architectural one; that the vending-machine
benchmark is answered by one $\dok$; and that the present culture's
hedges live in its harnesses.

Not claimed: that Lean-style verification covers only sequential or
linear-time software --- it does not, and Section~\ref{inh:sec:verification}
opens by saying so; that the alignment of the ladder with the
linear/branching/causal spectrum is a theorem beyond its first two
rungs; that Observation~\ref{inh:obs:rsi} is a theorem, though the
materials for one are now in the book; or that any transformer
\emph{cannot} be made to inhabit a race --- only that inhabiting is a
property of the witness language and not of the model's size.
The thing this book can do about the saddle point is what the Turn
did: build the witness language in which a hypothesis can be a race, a
certificate can run, and a hedge can be inherited, and show that it
runs on a machine.
Whether anyone leaves the saddle point is not a question a book
answers.

% ===================================================================
